\documentclass[10pt]{article}
\usepackage[margin=0.75in]{geometry}
\usepackage{times}
\usepackage{enumitem}
\usepackage{titlesec}
\usepackage{hyperref}
\usepackage{parskip}
\setlength{\parskip}{3pt}
\setlist{leftmargin=1.3em, itemsep=0pt, topsep=1pt}
\titlespacing*{\section}{0pt}{7pt}{3pt}
\titleformat{\section}{\normalfont\large\bfseries}{}{0pt}{}
\pagestyle{empty}

\title{\vspace{-1.6em}\textbf{Two Extensions of Reusable Clinical Skill Libraries:\\
Weak-to-Strong Transfer and Validated Skill Updates}}
\author{\vspace{-0.4em}Yiming Xiong \quad Jinkun Yang}
\date{}

\begin{document}
\maketitle
\vspace{-1.1em}

\section*{Problem and why it matters}
Clinical LLM agents diagnose patients and, between encounters, revise a persistent
\emph{skill library} instead of their weights. \textsc{DxSkill}, the first phase of this
project (complete, submitted to ACL ARR), built this framework: a diagnosing agent
reflects on each MIMIC-CDM patient encounter and revises a persistent, human-readable
skill library under five feedback conditions (from trajectory-only reflection to
clinical guidelines and personalized teacher advice), with model weights kept fixed
throughout. We found that accuracy depends heavily on which feedback is used
(77.5\%$\to$86.9--89.1\%) and that skill revision is not reliably safe -- a useful skill
can be silently deleted, and a transient tool error can harden into a permanent
restriction, even with outcome feedback available. This proposal
covers the next phase: two independent directions, (A) weak-to-strong skill transfer and
(B) buffer-validated skill updates, described below.

\textbf{Direction A} (inspired by OpenAI's 2023 weak-to-strong generalization work): can a
weaker model teach a stronger diagnostic agent by authoring and maintaining the skill
library the strong model executes, rather than supervising its predictions directly?
\textbf{Direction B}: can skill revision be made safer by validating each proposed update
on a holdout buffer before committing it?

\section*{Data, and whether we have access}
MIMIC-CDM 1.1 (from MIMIC-IV): 250 training / 400 held-out patients across appendicitis,
cholecystitis, diverticulitis, and pancreatitis.

\section*{Approach and baselines}
We define ``weak'' and ``strong'' by each model's measured no-skill diagnostic accuracy on
MIMIC-CDM, not by parameter count: \textbf{weak} = Qwen3-30B-A3B-Instruct-2507-FP8
(77.5\%); \textbf{strong} = Qwen3.8-27B-FP8 with thinking enabled (85.1\%). Our baselines
are DxSkill's own feedback conditions (up to 86.9\% with student self-revision) and
DxEvolve (Ren et al., 2026), an existing experience-distillation method we already
compare against.

\textbf{A -- weak-to-strong skill transfer (primary direction).} We test every
combination of a weak or strong model as the diagnosing agent with a weak or strong model
as the skill reviser -- four combinations in total -- holding feedback (reference
diagnosis) fixed. The case where a weak model revises the skill library for a strong
diagnosing agent is the key untested condition.

\textbf{B -- validation-buffer gating (secondary direction).} Before committing each
proposed skill update, we evaluate the candidate library on a fixed holdout set of 800
cases, kept separate from both the training stream and the 400-patient evaluation set. An
update is kept only if accuracy does not drop by more than a tolerance $\varepsilon$,
otherwise we roll back to the previous version of the library.

\section*{How we'll evaluate it}
\textbf{A:} accuracy on the frozen 400-patient set for each of the four
diagnosing-model/skill-reviser combinations, with paired patient-bootstrap intervals (as
in our ACL ARR submission), to locate where weak-to-strong transfer helps or fails.
\textbf{B:} accuracy on a fixed holdout set of 800 cases.

\section*{Milestones through Nov 30}
\begin{itemize}
\item \textbf{Oct 20--Nov 2 (A):} Run all four diagnosing-model/skill-reviser combinations.
\item \textbf{Nov 3--16 (A):} Analyze whether gains trace to execution or reflection strength.
\item \textbf{Nov 10--23 (B):} Implement the holdout buffer and checkpointed rollback; run a tolerance ($\varepsilon$) sweep.
\item \textbf{Nov 24--30:} Consolidate whichever results are ready into a progress report; Direction B may continue past Nov 30.
\end{itemize}

\section*{What we need from the course staff}
GPU resources on the Miata server are enough for both directions; no additional compute
or API budget is needed.

\end{document}
